%!TEX root = ../../../standalone/applications/app-feelpp-toolbox.tex
\ApplicationFicheHeading{app-feelpp-toolbox}
\textbf{app-feelpp-toolbox: Multi-Physics Toolbox Benchmark}
\label{sec:app:specs:app-feelpp-toolbox}

% Generated by scripts/sync_applications.py; edit the CSV snapshot.
\begin{FicheMetadataTable}
id & app-feelpp-toolbox \\
name & Test Feel++ toolboxes at scale \\
Partners & Unistra \\
PC & PC1 - ExaMA \\
Responsible (Permanent) & V. Chabannes; C. Prud'homme \\
WP7 Engineer & Javier Cladellas (UNISTRA) \\
work\_package & WP1, WP3 \\
application\_type & demonstrator \\
purpose & test heat(eye2brain,thermalbridge), elasticity(nafems LE-10), thermoelectric(hifimagnet) and fluid (fda) \\
Method-Algorithm WP1 & unstructured mesh, mesh adaptation, multiphysics coupling, finite element \\
Method-Algorithm WP3 & domain decomposition methods, multiphysics coupling, algebraic multiphysics coupling, preconditioning \\
inputs & Gmsh, JSON workflow \\
outputs & HDF5 fields, JSON small-data \\
metrics & time-to-solution, benchmark-verification \\
status & benchmark-ready \\
Benchmark scope & multi-node \\
Framework & Feel++ \\
parallel\_framework & MPI \\
notes & HiFiMagnet thermo-electric work is detailed in app-feelpp-mi300a-fem, app-feelpp-hifimagnet-inverse and app-feelpp-hifimagnet-optcontrol. \\
\end{FicheMetadataTable}


\subsection{Description of the benchmark}
\textbf{Catalogue status:} Benchmark-ready | \textbf{Partners:} Unistra | \textbf{Responsible:} V. Chabannes; C. Prud'homme

This is an umbrella application for the installed \Feelpp toolboxes.
The catalogue status records operational readiness; validation and performance
results must be assessed separately for each physics case.

\textbf{Purpose:} Test \Feelpp toolboxes at scale across multiple physics domains (demonstrator).

\paragraph{Available cases and results.}
The application repository provides JSON configurations for the cases below.
Each runs an installed Feel++ toolbox executable. The results available for
each case are listed separately from its launch configuration.

\begin{table}[!ht]
\centering\small
\begin{tabularx}{\textwidth}{@{}>{\raggedright\arraybackslash}p{0.23\textwidth}>{\raggedright\arraybackslash}p{0.32\textwidth}>{\raggedright\arraybackslash}X@{}}
\toprule
\textbf{Physics and case} & \textbf{Executable and launch configuration} & \textbf{Available results and related fiche} \\
\midrule
Heat: thermal bridges, ISO 10211 &
\nolinkurl{feelpp_toolbox_heat}\newline
\href{https://github.com/numpex/apps-feelpp/blob/main/src/app-feelpp-discr-1/ThermalBridgesENISO10211/heat_thermalbridges.json}{JSON configuration} &
\href{https://numpex.github.io/apps-feelpp/apps-feelpp/leaves/gaya-feelpp_toolbox_heat-ThermalBridgesENISO10211-2026_01_16T14_41_50/leaf.html}{Gaya report, 16 January 2026};
\Cref{sec:app:specs:app-feelpp-discr-1}. \\
\addlinespace
Heat--fluid: Eye2Brain &
\nolinkurl{feelpp_toolbox_heatfluid}\newline
\href{https://github.com/numpex/apps-feelpp/blob/main/src/app-feelpp-toolbox/heatfluid/eye2brain/heatfluid_eye2brain.json}{JSON configuration} &
Coupled heat transfer and flow. A configuration and plot definition exist; no matching report is currently exposed by the production dashboard. \\
\addlinespace
Solid mechanics: NAFEMS LE10 &
\nolinkurl{feelpp_toolbox_solid}\newline
\href{https://github.com/numpex/apps-feelpp/blob/main/src/app-feelpp-toolbox/solid/nafemsle10/elasticity_nafemsle10.json}{JSON configuration} &
Linear-elasticity case; a dedicated validation tolerance and archived campaign data remain to be attached to this fiche. \\
\addlinespace
Thermoelectric: HiFiMagnet HL-31 &
\nolinkurl{feelpp_toolbox_thermoelectric}\newline
\href{https://github.com/numpex/apps-feelpp/blob/main/src/app-feelpp-toolbox/thermoelectric/hifimagnet/thermoelectric_hl31.json}{JSON configuration} &
Historical CPU results in \Cref{sec:app:specs:thermoelectric:hl-31}; the MI300A campaign is a separate planned application. \\
\bottomrule
\end{tabularx}
\caption{Toolbox configurations and available results.}
\label{tab:app-feelpp-toolbox:cases}
\end{table}

The FDA nozzle belongs to the specification in
\Cref{sec:app:specs:app-feelpp-discr-2}. Its launch configuration and
validation are still pending.

\textbf{Key Methods:}
\begin{itemize}
\item WP1: Multi-physics coupling, finite elements
\item WP3: Domain decomposition, iterative solvers
\end{itemize}

\textbf{Framework:} Feel++, PETSc. The cases use MPI parallelism and
case-specific mesh partitions and solver options.

\textbf{Note:} This demonstrator consolidates multiple physics applications to provide a comprehensive evaluation of the \Feelpp toolbox ecosystem at HPC scale.

\subsection{Benchmarking methodology}
Each case supplies an executable, input-file dependencies, native or Apptainer
execution settings, resource/mesh/order sweeps, and report stages to
\textit{feelpp.benchmarking}. Toolbox timers distinguish initialization,
assembly/solve and post-processing; the selected stages differ by physics.
Verification outputs include physical measures where the configuration
requests them. A configured parameter sweep does not establish that every point
completed successfully.

For each campaign, record the application and Feel++ versions, input
dataset identifiers, solver tolerances, machine/process layout, container
digest or native toolchain, repeated-run variability, and raw reports.
WP1 supplies the physical/discretization verification; WP3 assesses solver
and parallel performance; WP7 provides the execution/reporting workflow.

\subsection{Input/Output Dataset Description}

\subsubsection{Input Data}
Gmsh geometries or partitioned JSON/HDF5 meshes, toolbox CFG/JSON model and
solver files, and the case-specific workflow and plotting configurations.
The application repository refers to external datasets; it does not bundle
all meshes or the toolbox executables.

\subsubsection{Output Data}
Visualization fields and case-dependent CSV measures, timing tables,
JSON metadata and generated reports. Output formats and measured phases are
declared by each launch configuration.

\subsection{Results summary}
The thermal-bridges case has the linked public Gaya report. The HL-31 fiche
retains its existing CPU campaign; neither that result nor a working launch
configuration establishes completed GPU or full-system scaling.
The \href{https://numpex.github.io/apps-feelpp/apps-feelpp/index.html}{production dashboard}
currently exposes four reports dated December 2025--January 2026: thermal
bridges with the heat toolbox, plus three I/O runs. Its deployment date is
not a benchmark run date. This umbrella fiche does not introduce a separate
aggregate performance result or duplicate the component benchmarks.

% Generated by scripts/feelpp_catalogue.py; edit data/feelpp-mesh-catalogue.json.

\subsection{Application package and shared meshes}
\begin{FicheMetadataTable}
Repository entry & \nolinkurl{apps/app-feelpp-toolbox/} (local review layout in apps-feelpp; publication pending) \\
Delivery model & toolbox config \\
Available material & Four physics configurations indexed separately; public heat report available. \\
Results status & Results from an earlier campaign. \\
\hyperref[mesh:feelpp:thermal-bridges]{ISO thermal bridges} & ISO 10211 case geometry; preserve material interfaces and boundary labels. Archive geometry/model version, length unit, partition counts, cell/DoF statistics and ISO observables. \\
\hyperref[mesh:feelpp:eye2brain]{Eye2Brain} & Ocular geometry; physical boundary and subdomain labels. Record units, common mesh identity across apps, partition recipes, model revision and per-order DoFs. \\
\hyperref[mesh:feelpp:nafems-le10]{NAFEMS LE10} & Solid mechanics reference geometry and boundary markers. Pin geometry/model, load/constraint masks, stress observation point and validation tolerance. \\
\hyperref[mesh:feelpp:hl31]{HiFiMagnet HL-31} & Magnet conductor/cooling geometry; SI units in shared forward model. Preserve conductor/cooling/contact/electrode markers, geometry scale, material laws and partition recipes. \\
\end{FicheMetadataTable}
\paragraph{Next contribution.} Attach case-specific validation observables and thresholds and resolve external model/mesh dependencies. Link component results rather than duplicate campaigns.

